% problem-000183 5 氢能源汽车热力学
\section*{5. 氢能源汽车热力学}
碱⼟或稀⼟元素(A)和过渡⾦属(B)可以形成多种合⾦，⼴泛应⽤于催化、储氢等领域。下图给出某合⾦的理想结构沿不同⽅向的投影⽰意图，此结构属六⽅晶系，晶胞参数 $\displaystyle a = 5 4 0 . 9 \mathrm { p m } , \ c = 4 3 0 . 0 \mathrm { p m }$ 。其中，⼤球为A原⼦，⼩球为B原⼦，圆圈表⽰空位。

（图：）
(a)

（图：）
(b)

（图：）
(c)

（图：）
(d)
图5.1 某合⾦理想结构沿不同⽅向的投影⽰意图（其中，⼤球为A原⼦，⼩球为B原⼦，圆圈表⽰空位。）

(a) 晶胞；(b) 沿�⽅向投影；(c) A 和 B 混合排列层 $\left( \mathrm { L } _ { 1 } | \frac { \mathbf { \equiv } } { \mathbf { z } } \right)$ ；(d) 全部由B原⼦组成的层 $\left( \mathrm { L } _ { 2 } \overleftrightarrow { \underline { { \mathbf { \Lambda } } } } \right)$ 5-1 写出该合⾦的组成。

\subsection*{5-2}
⼰知原⼦A处在晶胞原点，写出晶胞中所有B原⼦的坐标参数。（排序要求：先按�从⼩到⼤；�相同时，按照�从⼩到⼤；�相同时，按照�从⼩到⼤。)

\subsection*{5-3}
计算同⼀层内B原⼦的最短距离和相邻层间B 原⼦的最短距离（单位：pm）。

\subsection*{5-4}
该合⾦有良好的储氢性能。研究发现，氢原⼦(H)占据如下位置（结构中实际位置与此有所偏离，但计量关系⼀致）； $\mathrm { L _ { 1 } }$ 层中所有菱形的中⼼ $( \mathrm { i } \mathrm { \vec { E } \mathcal { H } H _ { R } } )$ $\mathrm { L } _ { 2 }$ 层中的三⻆形中⼼且只占⼀半 $( \mathrm { i } \mathrm { \vec { E } \partial \mathrm { \partial \cdot \partial \mathrm { J } H _ { T } ) } }$ o

\subsection*{5-4}
-1 写出晶胞中两种氢原⼦ $\mathrm { H } _ { \mathrm { R } }$ 和 $\mathrm { H } _ { \mathrm { T } }$ 的数⽬。

\subsection*{5-4}
-2 将氘化（即⽤D取代H）的样品进⾏中⼦衍射，发现 $\mathrm { L _ { 1 } }$ 层中⾦属原⼦和氢原⼦位置与上述结构吻合；⽽$\mathrm { L } _ { 2 }$ 层中氢原⼦的排布完全有序：倘若某⼀层中D占据全部顶点朝上的三⻆形中⼼（记为 $\mathrm { L } _ { 2 \Delta } )$ ），则其邻近的 $\mathrm { L } _ { 2 }$ 层中D占据全部顶点朝下的三⻆形中⼼（记为 $\operatorname { L } _ { 2 \nabla } )$ ）。晶体结构按照 $\therefore \mathrm { L _ { 1 } L _ { 2 \Delta } L _ { 1 } L _ { 2 \nabla } } . .$ 进⾏周期性排列。写出该氘化物的晶胞参数，假设D代不影响A和B原⼦的位置。
